\textbf{COLUNI 2022 -- Questão 29}\par\medskip O Brasil é reconhecido por especialistas como o País de maior biodiversidade do planeta. Porém, essa enorme biodiversidade vem sendo ameaçada de forma crescente. Nosso País apresenta atualmente um Sistema Nacional de Unidades de Conservação considerável: são mais de 1,4 milhão de km$^2$, somadas apenas as Unidades de Conservação (UC) federais e estaduais. Entretanto, 27\% desse total referem-se a áreas de proteção ambiental, que, apesar do nome, em sua maioria não garantem uma proteção adequada aos ecossistemas e seus componentes. Além disso, a distribuição das UC pelos seis biomas continentais e pela zona costeira e marinha é bastante desigual. As UC federais e estaduais de proteção integral (destinadas à manutenção dos ecossistemas livres de alterações causadas por interferência humana) correspondem a 9,8\% da Amazônia, mas a apenas 1,0\% da Caatinga, 2,9\% do Cerrado, 2,1\% da Mata Atlântica, 1,1\% do Pampa, 2,3\% do Pantanal e 0,1\% da Zona Costeira e Marinha, que engloba o mar territorial e a zona econômica exclusiva.

\par

ROMA, J. C. e VIANA, J. P. Repositório do Conhecimento do IPEA: Conservação desbalanceada entre os biomas. 2009. Disponível em: http://repositorio.ipea.gov.br/handle/11058/8371. Acesso em: 01 set. 2021.

\par

Sobre a ecologia e a conservação dos biomas brasileiros, assinale a afirmativa INCORRETA:\begin{enumerate}
\item A Mata Atlântica é um dos biomas brasileiros mais impactados pelas ações humanas e seu histórico de
      degradação está intimamente associado ao desmatamento e à expansão da agricultura e de áreas
      urbanas.
\item A Floresta Amazônica, maior floresta tropical do mundo, caracteriza-se pela ocorrência de diferentes
      ecossistemas, como as matas de igapó, matas de várzea e matas de terra firme, em grande parte,
      determinadas pelo relevo e pela hidrografia.
\item O Cerrado ocupa extensa área da região central do Brasil e apresenta espécies bem adaptadas a altas
      intensidades luminosas e duas estações bem definidas. A exploração de carvão e a mineração são
      atividades de grande impacto ambiental para esse bioma.
\item O Pantanal está localizado na bacia do Rio Paraguai e as diversas paisagens formadas nessa região
      apresentam similaridades com outros biomas, como a Amazônia e o Cerrado. O desmatamento e o fogo
      são os principais impactos da ação humana nesse bioma.\end{enumerate}
